\documentclass[12pt]{article}
\usepackage{amsmath}
\usepackage{geometry}
\geometry{letterpaper, margin=1in}
\usepackage{fancyhdr}
\usepackage{setspace}
\usepackage{hyperref}

\pagestyle{fancy}
\fancyhf{}
\rhead{33-120 Science and Science Fiction}
\lhead{Spacetime Team Project (25 points)}
\setlength{\headheight}{15pt}

\begin{document}

\vspace{1cm}
{\large{\noindent Demonstrating the Equivalence of a Gravitational Field and an
Accelerating Reference Frame}}

\noindent\rule{6.5in}{0.4pt}

\vspace{0.5cm}
{\bf{\noindent Finding Your Team Assignment \& Coordinating with Your Team to
Pick Up Equipment}}
\begin{itemize}
    \item You should have groups assigned by Canvas.
    \item Contact your team members and designate someone to pick up your packet
    of equipment.
    \item Pick up equipment after class on Monday, Sept 16, or in Prof
    Breiviks's office, Wean 8319 during office hours from 2pm-3pm.
    \item You may do the project at your convenience, either during the week
    before the designated project day or on the designated days (Sept 18,20)
    during regular class time.
\end{itemize}

{\bf{\noindent Making and Recording the Measurements}}
\begin{itemize}
    \item All measurements should be done during our regularly scheduled class
    time unless agreed otherwise with your team.
    \item Decide as a team which TWO elevators you will use for measurements.
    You may use any elevators in academic buildings (but not in administrative
    buildings or the University Center).
    \item Provide enough information about the specific location of your
    elevators (e.g., directions, photos, room numbers, etc.).
    \item Check the “zero” on the force scale:
    \begin{itemize}
        \item The larger hook is at the top end, fixed inside the scale, and the
        smaller hook is attached to the spring.
        \item Hold the scale freely and adjust the zero if necessary by
        loosening the locking ring.
    \end{itemize}
    \item Understand the scale readings:
    \begin{itemize}
        \item The maximum force is 1.0 N.
        \item Fine scale divisions are either 0.01 N or 0.02 N.
    \end{itemize}
    \item Measure the gravitational force on a 50-gram mass at rest.
    \item Measure the apparent force on the object in two non-inertial reference
    frames (elevators):
    \begin{itemize}
        \item Take the mass and scale on a ride in the elevator.
        \item Measure the force at the instant the elevator starts to move up.
        \item Measure the force when the elevator comes to a stop while moving
        down.
        \item Only record readings that are greater than the reading at rest.
    \end{itemize}
\end{itemize}


{\bf{\noindent Comparing the Results}}

\noindent Answer the questions based on your measurements. Do you detect any
significant difference between the Starting/UP and Stopping/DOWN measurements?
Are the two elevators significantly different from each other?

\vspace{0.5cm}
{\bf{\noindent Submitting the Results}}

\noindent Turn in one data sheet per team via file upload to Canvas on the
scheduled due date. Return your equipment to Prof Breivik in class on Monday,
Sept 23, or during office hours on Tuesday Sept 24. 

\newpage

\section*{\noindent Force Measurements in Accelerating Reference Frames}

\textbf{(Fill in this data sheet, and turn in the results on Canvas.)}

\vspace{0.5cm}
\noindent \textbf{Names of team members} (print clearly):

\vspace{0.5cm}
\noindent \textbf{Team Number}: \_\_\_\_

\vspace{0.5cm}
\noindent \textbf{Force on 50 gram object at rest} = \_\_\_\_ N

\begin{tabbing}
\textbf{Elevator 1 location} \hspace{2.5cm} \= \textbf{Elevator
2 location} \\
Building: \_\_\_\_\_\_\_\_\_\_\_\_\_\_\_\_\_\_ \hspace{2cm} \= Building: \_\_\_\_\_\_\_\_\_\_\_\_\_\_\_\_\_\_ \\
Specific Location: \_\_\_\_\_\_\_\_\_\_\_\_\_\_\_\_\_\_ \hspace{2cm} \= Specific Location: \_\_\_\_\_\_\_\_\_\_\_\_\_\_\_\_\_\_
\end{tabbing}

\vspace{0.5cm}

\begin{tabbing}
\textbf{Starting (Going Up)} \hspace{2cm} \= \textbf{Stopping
(Going Down)} \\
1. \_\_\_\_ N \hspace{4cm} \= 1. \_\_\_\_ N \\
2. \_\_\_\_ N \hspace{4cm} \= 2. \_\_\_\_ N \\
3. \_\_\_\_ N \hspace{4cm} \= 3. \_\_\_\_ N \\
4. \_\_\_\_ N \hspace{4cm} \= 4. \_\_\_\_ N \\
5. \_\_\_\_ N \hspace{4cm} \= 5. \_\_\_\_ N \\
\end{tabbing}

\vspace{0.5cm}

\begin{tabbing}
    \textbf{Starting (Going Up)} \hspace{2cm} \= \textbf{Stopping
    (Going Down)} \\
    1. \_\_\_\_ N \hspace{4cm} \= 1. \_\_\_\_ N \\
    2. \_\_\_\_ N \hspace{4cm} \= 2. \_\_\_\_ N \\
    3. \_\_\_\_ N \hspace{4cm} \= 3. \_\_\_\_ N \\
    4. \_\_\_\_ N \hspace{4cm} \= 4. \_\_\_\_ N \\
    5. \_\_\_\_ N \hspace{4cm} \= 5. \_\_\_\_ N \\
    \end{tabbing}
    
    \vspace{0.5cm}

\noindent \textbf{Questions for thought...}
\begin{itemize}
    \item For a given elevator, do you detect any significant difference between
    the Starting/UP and Stopping/DOWN measurements?
    \vspace{1.5cm}
    \item Are the two elevators that you tested significantly different from
    each other?
\end{itemize}

\end{document}